\section{Experimental}
\subsection{The Hydrolysis of a Polyimide}
Poly..... (Kapton) film (0.2g, thickness of x) was cut into pieces of approx. 3mm$^2$, these were mixed with a solution of NaOH () in a reaction tube. This tube was placed into a monowave reactor that was set to ramp for 20 minutes and hold for 20 minutes at 100 $\degree$C. A sample was taken and the remaining mixture was reacted again on the monowave for ramp 20 mins and hold 20 minutes at 100$\degree$C. There was no observable difference in the Kapton film. Then this was thoroughly washed with water and IR was collected.

\subsection{Cross-linking with TEAB and Ethylene carbonate}
A mixture of polyamic acid (0.51mL, conc) was mixed with ethylene carbonate () and tetraethylammonium bromide (). This was mixed in a monowave reactor set to ramp for 20 minutes and hold for 1 hour at 150 $\degree$C. This gave an orange suspension that was dried to give a pale orange powder as the final product (grams, yield).
IR (solid): 

\subsection{Cross-linking with $ZnCl_2$ and Ethylene carbonate}
A mixture of polyamic acid (0.51mL, conc) was mixed with ethylene carbonate () and Zinc chloride (). This was mixed in a monowave reactor set to ramp for 20 minutes and hold for 1 hour at 100 $\degree$C. This gave dark orange clear liquid that was dried to give a () as the final product (grams, yield).
IR (solid): 

\subsection{Cross-linking with $FeCl_3$ and Ethylene carbonate}
A mixture of polyamic acid (0.51mL, conc) was mixed with ethylene carbonate () and Iron(III) chloride (). This was mixed in a monowave reactor set to ramp for 20 minutes and hold for 1 hour at 100 $\degree$C. This gave a dark orange clear gel that was dried to give a () as the final product (grams, yield).
IR (solid): 